% p99-concurrencia.tex — Figura obligatoria 1 (R-30): P99 contra concurrencia.
%
% R-30.1 (decisión de Alix, al aprobar F7-e): la figura muestra LAS DOS VISTAS, no
% una nota al pie. Con c=50 en los perfiles B y C, hasta el 91 % de las consultas
% de la ventana son la primera consulta en frío de cada cliente (F7.3), así que
% parte de la subida de la curva es arranque en frío y no latencia del motor.
%   - línea continua: el P99 de la columna del CSV (lo que mide el manual);
%   - línea punteada: el P99 recalculado sin la primera consulta de cada cliente,
%     desde results/raw/, que es la vista en régimen.
% Donde el .dat trae `nan` no hay vista en régimen porque TODAS las consultas de
% la ventana eran primeras consultas: la línea punteada se interrumpe ahí
% (unbounded coords=jump, en el preámbulo).
%
% Los cuatro paneles van escritos uno a uno, SIN un macro con parámetros: dentro
% de un \newcommand, los '#' del cuerpo chocan con el analizador de claves de
% pgfplots ("Illegal parameter number in definition of \pgfkeys@..."). Es más
% largo, pero no depende de cómo pgfkeys reinterprete el cuerpo.

\begin{figure}[htbp]
  \centering
  \begin{tikzpicture}
    \begin{axis}[
        title={MongoDB · Q2 · perfil A}, title style={font=\small},
        xlabel={Concurrencia (clientes)}, ylabel={P99 (ms)},
        xmode=log, log basis x=10, xtick={1,10,50}, xticklabels={1,10,50},
        ymin=0, width=0.47\linewidth, height=5cm,
        legend pos=north west, legend cell align=left,
      ]
      \addplot[color=blue!60!black, mark=*, thick] table [x=concurrencia, y=p99_csv] {datos/p99_mongodb_Q2_A.dat};
      \addplot[color=blue!60!black, mark=o, densely dotted, thick]
        table [x=concurrencia, y=p99_sin_frias] {datos/p99_mongodb_Q2_A.dat};
      \legend{P99 del CSV, Sin la 1.\textsuperscript{a} en frío}
    \end{axis}
  \end{tikzpicture}\hfill
  \begin{tikzpicture}
    \begin{axis}[
        title={PostgreSQL · Q2 · perfil A}, title style={font=\small},
        xlabel={Concurrencia (clientes)}, ylabel={P99 (ms)},
        xmode=log, log basis x=10, xtick={1,10,50}, xticklabels={1,10,50},
        ymin=0, width=0.47\linewidth, height=5cm,
        legend pos=north west, legend cell align=left,
      ]
      \addplot[color=red!60!black, mark=*, thick] table [x=concurrencia, y=p99_csv] {datos/p99_postgresql_Q2_A.dat};
      \addplot[color=red!60!black, mark=o, densely dotted, thick]
        table [x=concurrencia, y=p99_sin_frias] {datos/p99_postgresql_Q2_A.dat};
    \end{axis}
  \end{tikzpicture}

  \begin{tikzpicture}
    \begin{axis}[
        title={MongoDB · Q2 · perfil C}, title style={font=\small},
        xlabel={Concurrencia (clientes)}, ylabel={P99 (ms)},
        xmode=log, log basis x=10, xtick={1,10,50}, xticklabels={1,10,50},
        ymin=0, width=0.47\linewidth, height=5cm,
        legend pos=north west, legend cell align=left,
      ]
      \addplot[color=blue!60!black, mark=*, thick] table [x=concurrencia, y=p99_csv] {datos/p99_mongodb_Q2_C.dat};
      \addplot[color=blue!60!black, mark=o, densely dotted, thick]
        table [x=concurrencia, y=p99_sin_frias] {datos/p99_mongodb_Q2_C.dat};
    \end{axis}
  \end{tikzpicture}\hfill
  \begin{tikzpicture}
    \begin{axis}[
        title={PostgreSQL · Q2 · perfil C}, title style={font=\small},
        xlabel={Concurrencia (clientes)}, ylabel={P99 (ms)},
        xmode=log, log basis x=10, xtick={1,10,50}, xticklabels={1,10,50},
        ymin=0, width=0.47\linewidth, height=5cm,
        legend pos=north west, legend cell align=left,
      ]
      \addplot[color=red!60!black, mark=*, thick] table [x=concurrencia, y=p99_csv] {datos/p99_postgresql_Q2_C.dat};
      \addplot[color=red!60!black, mark=o, densely dotted, thick]
        table [x=concurrencia, y=p99_sin_frias] {datos/p99_postgresql_Q2_C.dat};
    \end{axis}
  \end{tikzpicture}

  \caption[P99 contra concurrencia, en dos vistas]{%
    \textbf{P99 contra concurrencia, en dos vistas (R-30.1).} La línea continua es
    el P99 tal como lo mide el manual (columna del CSV). La punteada es el mismo
    P99 \textbf{sin la primera consulta de cada cliente}, recalculado desde
    \texttt{results/raw/}. La separación entre las dos líneas es el arranque en
    frío: el calentamiento de 60~s del manual no alcanza con \mbox{c=50} en los
    perfiles B y C, donde hasta el 91~\% de la ventana son primeras consultas
    (F7.3). En MongoDB · perfil C · \mbox{c=50} la línea punteada se interrumpe
    porque \emph{todas} las consultas de la ventana eran primeras consultas y no
    existe un P99 en régimen que dibujar. \textbf{Leer solo la línea continua
    sobreestima la latencia del motor en esos puntos.} Datos:
    \texttt{report/datos/p99\_*.dat}.}
  \label{fig:p99-concurrencia}
\end{figure}
